\section{Selecting a generator on the boundary ray}
\label{sec:generator}

For a field \(K\) containing \(\C\), a real-valued valuation is a map
\(v:K^\times\to\R\) satisfying
\[
v(fg)=v(f)+v(g),\qquad
v(f+g)\geq\min\{v(f),v(g)\}\quad(f+g\ne0).
\]
All valuations in this section are trivial on \(\C^\times\). The zero
valuation is allowed. If a homogeneous polynomial has the form
\(f x^a y^b\), we apply \(v\) to its coefficient \(f\); the formal
variables \(x,y\) record degrees only.

The relevant boundary ray consists of degrees \((0,b)\). A sequence
approaches that ray after normalization when its first degree is
negligible compared with its total degree.

\begin{lemma}
\label{lem:endpoint}
Let \(R\subset K[x,y]\) be a finitely generated \(\C\)-subalgebra,
graded by \(\N^2\) with \(\deg x=(1,0)\), \(\deg y=(0,1)\), and
\(R_{0,0}=\C\). Let \(v\) be a real-valued valuation of \(K\), trivial
on \(\C^\times\). Suppose that nonzero homogeneous elements
\[
s_j=f_jx^{a_j}y^{b_j}\in R,\qquad d_j=a_j+b_j>0,
\]
satisfy
\[
\frac{a_j}{d_j}\longrightarrow0,
\qquad
\frac{v(f_j)}{d_j}\longrightarrow\gamma\in\R.
\]
Then every finite set of nonzero homogeneous generators of \(R\)
contains a generator \(f x^0y^b\), with \(b>0\), such that
\[
\frac{v(f)}b\leq\gamma.
\]
\end{lemma}

\begin{proof}
Discard constant generators and write the others as
\[
r_i=g_ix^{u_i}y^{w_i},
\qquad u_i,w_i\in\N,\qquad u_i+w_i>0.
\]
There are finitely many indices \(i\). Because the generators are
homogeneous, \(s_j\) can be expressed as a sum of monomials of degree
\((a_j,b_j)\). One nonzero monomial has coefficient valuation at most
\(v(f_j)\): otherwise the valuation inequality for the sum would
contradict its coefficient \(f_j\). Let \(e_{ij}\) be the exponents
of the generators in such a monomial. Nonzero scalar coefficients
have valuation zero, so
\begin{equation}
\label{eq:monomial}
\sum_i e_{ij}u_i=a_j,\qquad
\sum_i e_{ij}w_i=b_j,\qquad
\sum_i e_{ij}v(g_i)\leq v(f_j).
\end{equation}

We compare the factors with \(u_i>0\) and those with \(u_i=0\).
Every positive \(u_i\) is at least one; hence
\[
\sum_{u_i>0}e_{ij}\leq a_j.
\]
Since the generating set is finite, there are constants \(C_1,C_2\geq0\)
such that these factors contribute at most \(C_1a_j\) to the second
degree and at least \(-C_2a_j\) to the valuation. For example, take
\[
C_1=\max_i w_i,\qquad
C_2=\max_i\max\{0,-v(g_i)\}.
\]
The set of nonconstant generators is nonempty because \(d_j>0\).

Some generator must have \(u_i=0\). Otherwise
\eqref{eq:monomial} would give \(b_j\leq C_1a_j\), contrary to
\(a_j/d_j\to0\). For every such generator \(w_i>0\), so we may define
\[
\lambda=\min_{u_i=0}\frac{v(g_i)}{w_i}.
\]
Let \(c_j\) be the second degree contributed by the factors with
\(u_i>0\). Thus \(0\leq c_j\leq C_1a_j\). The remaining factors
contribute second degree \(b_j-c_j\), and their valuation is at least
\(\lambda(b_j-c_j)\). Equation~\eqref{eq:monomial} therefore gives
\[
v(f_j)\geq
\lambda b_j-\bigl(C_2+|\lambda|C_1\bigr)a_j.
\]
Divide by \(d_j\) and pass to the limit. Since \(b_j/d_j\to1\), we
obtain \(\gamma\geq\lambda\). A generator realizing the finite minimum
\(\lambda\) proves the assertion.
\end{proof}

The proof needs the finite list of actual coefficients \(g_i\).
It does not follow from knowing only a limiting cone of numerical
classes. This is the feature that will matter for divisor section
rings in the next subsection.

\subsection{Removing a birational correction}
\label{sec:geometric}

We now specify the geometric data to which the lemma applies. The main
point is that the correction on the base need not have valuation zero.
It will disappear under birational pushforward before the last
valuation comparison.

All varieties below are integral projective varieties over \(\C\).
For a Cartier divisor \(B\) on a normal variety \(Y\), we identify
sections with rational functions:
\[
H^0(Y,\OO_Y(B))
=\{h\in\C(Y)^\times:\divisor_Y(h)+B\geq0\}\cup\{0\}.
\]
For two Cartier divisors \(B_1,B_0\), their bigraded section ring is
\begin{equation}
\label{eq:section-ring}
R(Y;B_1,B_0)
=\bigoplus_{a,b\in\N}
H^0\bigl(Y,\OO_Y(aB_1+bB_0)\bigr)x^ay^b
\subset\C(Y)[x,y].
\end{equation}
Projectivity and integrality give \(H^0(Y,\OO_Y)=\C\), as required
in Lemma~\ref{lem:endpoint}.

We also use rational powers of rational functions, interpreted as
elements of
\(\C(Y)^\times\otimes_{\Z}\Q\), written multiplicatively.
The divisor map and every valuation extend \(\Q\)-linearly to this
group. Each element becomes an ordinary rational function after a
positive integral power is taken. No extension of the function field
by actual roots is intended.

A valuation of \(\C(X)\) is centered on \(X\) if its valuation ring
dominates \(\OO_{X,\xi}\) for some point \(\xi\in X\). For a Cartier
divisor \(D\), define \(v(D)\) by applying \(v\) to a local defining
rational function at \(\xi\). This is independent of the equation,
because units have valuation zero, and extends \(\Q\)-linearly to
rational Cartier divisors. In particular, \(v(D)\geq0\) for every
effective rational Cartier divisor \(D\). The zero valuation is
centered at the generic point.

\begin{proposition}
\label{prop:geometric}
Let \(X,Y\) be smooth integral projective complex varieties, let \(W\)
be normal and projective, and let
\[
X\xleftarrow{\ p\ }W\xrightarrow{\ f\ }Y
\]
be morphisms, with \(p\) birational and \(f\) surjective and
equidimensional. Identify \(\C(Y)\) with its induced subfield of
\(\C(X)=\C(W)\). Fix a centered discrete valuation \(v\) of
\(\C(X)\), allowing also the zero valuation.

Let \(A\geq0\) be a rational divisor on \(Y\) such that
\(p_*f^*A=0\). Let \(D_0,D_1\) be rational Cartier divisors on \(X\)
and put
\[
D(t)=(1-t)D_0+tD_1\qquad(0\leq t\leq1).
\]
For each prime divisor \(Q\) of \(Y\), define
\begin{equation}
\label{eq:minimum-divisor}
P(t)=\sum_Q
\min_{E:\,f(E)=Q}
\frac{\coeff_E p^*D(t)}{\ord_E(f^*Q)}\,Q,
\end{equation}
where \(E\) ranges over the prime divisors of \(W\) dominating \(Q\).
Assume that \(P(t)=(1-t)P(0)+tP(1)\) throughout \([0,1]\).

Suppose that rational numbers \(t_j\in[0,1]\) tend to zero and that
\(h_j\in\C(Y)^\times\otimes_{\Z}\Q\) satisfy
\[
N_j=D(t_j)+\divisor_X(h_j)\geq0,\qquad v(N_j)\longrightarrow0.
\]
Finally, suppose that for some \(\epsilon\in\Q_{>0}\) and integer
\(k>0\), the divisors
\[
B_1=k\bigl(A+\epsilon P(1)\bigr),\qquad
B_0=k\bigl(A+\epsilon P(0)\bigr)
\]
are Cartier and the ring \(R(Y;B_1,B_0)\) is finitely generated.
Then there is an effective rational divisor
\[
N\sim_{\Q}D_0,\qquad v(N)=0.
\]
\end{proposition}

The minimum in \eqref{eq:minimum-divisor} is over a finite nonempty
set, since \(f\) is surjective and equidimensional. Only finitely many
base primes occur: the pullbacks of \(D_0,D_1\) have fixed finite
support. Each coefficient is a minimum of finitely many rational
affine functions, so \(P(t)\) is piecewise affine. In applications one
restricts to a sufficiently small rational segment ending at \(t=0\)
and rescales the parameter to obtain the stated affine hypothesis.
If \(Y\) is a point, the sum is empty and means \(P(t)=0\).

\begin{proof}
Write
\[
t_j=\frac{a_j}{d_j},\qquad d_j=a_j+b_j>0,
\qquad a_j,b_j\in\N.
\]
Scale the pair \((a_j,b_j)\) by a positive integer whenever necessary
to clear denominators. The ratio \(a_j/d_j=t_j\) is unchanged.
In particular we can arrange that
\[
s_j=h_j^{\epsilon k d_j}
\]
is an ordinary rational function on \(Y\).

The effective divisor \(N_j\) is rational Cartier because \(X\) is
smooth, so \(p^*N_j\geq0\). For a prime \(E\) of \(W\) dominating a
base prime \(Q\), let \(e_E=\ord_E(f^*Q)>0\). Taking its coefficient
in this pullback gives
\[
\frac{\coeff_E p^*D(t_j)}{e_E}
+\coeff_Q\divisor_Y(h_j)\geq0.
\]
Taking the minimum proves
\begin{equation}
\label{eq:base-effectivity}
P(t_j)+\divisor_Y(h_j)\geq0.
\end{equation}
The affine hypothesis now gives
\[
a_jB_1+b_jB_0=k d_j\bigl(A+\epsilon P(t_j)\bigr).
\]
Since \(A\geq0\), equation~\eqref{eq:base-effectivity} shows that
\(s_jx^{a_j}y^{b_j}\) is a nonzero homogeneous element of
\(R(Y;B_1,B_0)\).

The valuation of \(D(t)\) is affine in \(t\), since the two divisors
have fixed support and valuation is \(\Q\)-linear. Thus
\[
v(h_j)=v(N_j)-v(D(t_j))\longrightarrow-v(D_0),
\qquad
\frac{v(s_j)}{d_j}\longrightarrow-\epsilon k\,v(D_0).
\]
Apply Lemma~\ref{lem:endpoint} to this section ring and to the
restriction of \(v\) to \(\C(Y)\). The lemma only uses this restriction
as a real-valued valuation. We obtain a
generator represented by a rational function \(r\), of degree
\((0,b)\) with \(b>0\), such that
\[
\frac{v(r)}{\epsilon k b}\leq-v(D_0).
\]
Put \(h_*=r^{1/(\epsilon k b)}\) in the rational-power group.
The section inequality for \(r\) gives
\begin{equation}
\label{eq:endpoint-effectivity}
P(0)+A/\epsilon+\divisor_Y(h_*)\geq0.
\end{equation}

We have selected the endpoint section on \(Y\). It remains to show that
it gives an effective divisor on \(X\) without retaining the correction
\(A/\epsilon\).

A prime of \(W\) is horizontal if it dominates \(Y\), and vertical
otherwise. Principal divisors pulled back from \(Y\) have no horizontal
components. Consequently the horizontal coefficients of \(p^*D(t_j)\)
are nonnegative, by effectivity of \(p^*N_j\); the same is true for
\(p^*D_0\) by passage to the limit. Every vertical prime dominates a
base prime, by equidimensionality. For such a prime \(E\) over \(Q\),
the minimum defining \(P(0)\) and
\eqref{eq:endpoint-effectivity} show that its coefficient in
\[
p^*D_0+
f^*\bigl(\divisor_Y(h_*)+A/\epsilon\bigr)
\]
is nonnegative. This divisor is therefore effective.

Push it forward by \(p\). The assumption \(p_*f^*A=0\) removes the
correction, while pullback and pushforward of the principal divisor
give \(\divisor_X(h_*)\). Hence
\[
N=D_0+\divisor_X(h_*)\geq0,\qquad N\sim_{\Q}D_0.
\]
Centering gives \(v(N)\geq0\). The selected generator gives
\(v(N)=v(D_0)+v(h_*)\leq0\), proving equality.
\end{proof}

The proof never evaluates the correction \(f^*A\) at \(v\).
Its removal is a divisor identity under pushforward. This permits an
exceptional correction with positive valuation on a higher model.

